\begin{tikzpicture}[
  font=\footnotesize,
  cell/.style={draw=black!60, thick, minimum width=42mm, minimum height=7mm,
               align=center, inner sep=2pt, anchor=south west},
  canary/.style={cell, fill=probe!14, draw=probe!60!black},
  ret/.style={cell, fill=culprit!12, draw=culprit!65},
  buf/.style={cell, fill=black!5},
  ar/.style={-{Stealth[length=1.7mm]}, semithick, draw=black!70},
  lbl/.style={font=\scriptsize, text=black!65, align=left},
]
  % stack grows down: high address at the top
  \node[buf]    (buf) at (0,0)    {local buffer};
  \node[canary] (can) at (0,0.72) {stack canary};
  \node[ret]    (ret) at (0,1.44) {saved return address};
  \node[lbl, anchor=south] at (2.1,2.26) {high address};
  \node[lbl, anchor=north] at (2.1,-0.1) {low address};
  \draw[ar, draw=culprit!70] (-0.45,0.1) -- (-0.45,2.0)
     node[midway, left, lbl, text=culprit, rotate=90, anchor=south]{overflow writes up};
  \node[lbl, anchor=west, text=black!60, align=left] at (4.7,0.72)
     {a linear overflow reaches the return\\ address only by overwriting the canary,\\
      and the function checks the canary\\ before returning, so the change is caught};
\end{tikzpicture}
